% ---------------------------------------------------------------------------
%  fig5.tex -- Figure 5 of the CMOS chain paper (fig:fig3), the main results
%  figure: reliability, dissipation and the tradeoff between them.  Pure
%  vector.
%
%  Replaces PostReviewFigures/CompositeFigure3_v3.pdf, a matplotlib export
%  514.1 x 257.7pt shown at 6in, i.e. rescaled by 0.840, so its nominally 15pt
%  axis labels landed on the page at 12.6pt.  Authored here at exactly
%  \textwidth = 510pt: include with [width=\textwidth] in a figure*, no
%  rescaling.  pdfinfo reads 508.09 bp.
%
%  Layout follows the original: (a) top-left, (b) bottom-left sharing the L
%  axis, (c) full height on the right with its y axis on the right-hand side,
%  its top edge level with (a) and its bottom edge level with (b).
%
%  Colour: figvddA..D from figstyle.tex, the four viridis samples the
%  matplotlib figure used, so V_dd means the same colour in all three panels.
%
%  Data: fig5_{tau,dissipation,tradeoff,ref}.dat, written by prep_fig5.py from
%  ../PostReviewFigures/Figure 5 Data/{tau_err,dissipation}.txt.
%
%  Build: make fig5
% ---------------------------------------------------------------------------
\documentclass[border=0pt]{standalone}
% ---------------------------------------------------------------------------
%  figstyle.tex -- shared style for every figure in the CMOS chain paper.
%  \input this from each figN.tex. Do not restyle locally; change it here.
%
%  Design rule: every figure is authored at EXACTLY the width it occupies on
%  the page, so \includegraphics applies no scaling and a 9pt label in one
%  figure is a 9pt label in every other. The paper's geometry, measured from
%  revtex4-2 [reprint,aps,pre]:
%      \columnwidth = 246.0pt = 3.404in   -> \figcolwidth
%      \textwidth   = 510.0pt = 7.057in   -> \figfullwidth
%  Include single-column figures with [width=\columnwidth] and full-width
%  figures (figure*) with [width=\textwidth]. Nothing else.
%
%  Verifying the width: pdfinfo reports big points (72/in), not LaTeX pt
%  (72.27/in). A correct column figure reads 245.08 bp; a correct full-width
%  figure reads 508.09 bp. Those are the numbers to check, not 246 and 510.
%
%  Fonts: Computer Modern, i.e. the LaTeX default. Load no font package.
%  Class sizes for reference: body 10, \small 9, \footnotesize 8, \scriptsize 7.
% ---------------------------------------------------------------------------
\usepackage{tikz}
\usepackage{pgfplots}
\usepackage{amsmath}
\usepackage{braket}
\pgfplotsset{compat=1.17}
\usetikzlibrary{arrows.meta,positioning,calc,backgrounds,fit,decorations.pathreplacing}
\usepgfplotslibrary{colormaps}

% ---- the two canonical widths -------------------------------------------
\newlength{\figcolwidth}   \setlength{\figcolwidth}{246.0pt}
\newlength{\figfullwidth}  \setlength{\figfullwidth}{510.0pt}

% ---- force an exact output width ----------------------------------------
%  Put \figcanvas{<width>}{<height>} as the FIRST line inside a tikzpicture.
%  It pins the bounding box, so the emitted PDF is exactly that size no
%  matter what the content does, and no rescaling happens at include time.
\newcommand{\figcanvas}[2]{\useasboundingbox (0,0) rectangle (#1,#2);}

% ---- type sizes (true points, because scale is 1) -----------------------
\newcommand{\figaxislabel}{\small}       % 9pt  -- axis labels
\newcommand{\figticklabel}{\footnotesize}% 8pt  -- tick labels
\newcommand{\figannot}{\footnotesize}    % 8pt  -- in-figure annotation
\newcommand{\figpanel}{\small\bfseries}  % 9pt  -- panel letters (a), (b)

% ---- palette -------------------------------------------------------------
\definecolor{figblueA}{HTML}{1F4E9C}  % dark blue   -- emphasis curve 1
\definecolor{figblueB}{HTML}{3C8DBC}  % mid blue    -- emphasis curve 2
\definecolor{figgreen}{HTML}{9CC7A9}  % pale green  -- de-emphasised curves
\definecolor{figgray}{HTML}{7F7F7F}   % neutral rules and guides
\definecolor{figred}{HTML}{C0392B}    % the one warning/annotation colour
\definecolor{fignodeA}{HTML}{6E8B74}  % node dot, k=0 side
\definecolor{fignodeB}{HTML}{1F4E9C}  % node dot, k=1 side
\definecolor{figshadeA}{HTML}{EAEFEA} % k=0 panel background
\definecolor{figshadeB}{HTML}{DCE8F5} % k=1 panel background
% second operator colour: the reduced MPO (physical dimension d), as against
% the full MPO (physical dimension D) drawn in figblueA/figblueB.
\definecolor{figtealA}{HTML}{1B6E6A} % reduced-MPO outline
\definecolor{figtealB}{HTML}{9FCFCB} % reduced-MPO fill

% ---- the V_dd ramp -------------------------------------------------------
%  Four samples of viridis at 0.08, 0.36, 0.64 and 0.92, the ramp the
%  matplotlib composite used, so a given V_dd is the same colour in every
%  panel that shows one.  The values are read out of the content stream of
%  PostReviewFigures/CompositeFigure3_v3.pdf, so they are matplotlib's own.
\definecolor{figvddA}{rgb}{0.283197,0.115680,0.436115}  % V_dd/V_T = 1.1
\definecolor{figvddB}{rgb}{0.179019,0.433756,0.557430}  % V_dd/V_T = 1.2
\definecolor{figvddC}{rgb}{0.170948,0.694384,0.493803}  % V_dd/V_T = 1.3
\definecolor{figvddD}{rgb}{0.793760,0.880678,0.120005}  % V_dd/V_T = 1.4

% ---- line weights --------------------------------------------------------
\pgfplotsset{
  figframe/.style={line width=0.6pt},
  figcurve/.style={line width=1.0pt},
  figemph/.style={line width=1.4pt},
}
\tikzset{
  figpanellabel/.style={font=\small\bfseries, inner sep=1pt},
  figwire/.style={line width=0.9pt},
  figlead/.style={figgray, line width=0.5pt, densely dotted},
  figarrow/.style={-{Stealth[length=4pt,width=3pt]}, line width=0.7pt},
}

% ---- the common axis style ----------------------------------------------
%  Matches the tick conventions already used by the matplotlib figures:
%  inward ticks on all four sides, minor ticks visible, no grid.
\pgfplotsset{
  figaxis/.style={
    scale only axis,
    axis line style={line width=0.6pt},
    tick align=inside,
    tickpos=both,
    minor tick num=1,
    tick label style={font=\figticklabel},
    label style={font=\figaxislabel},
    legend style={font=\figannot, draw=none, fill=none},
    every axis plot/.append style={figcurve},
  },
  figdensity/.style={
    figaxis,
    % `viridis high res' is matplotlib's own 256-entry table (the plain
    % `viridis' key is an 18-point approximation to it), so a density panel
    % ported from matplotlib keeps the colours it had.
    colormap/viridis high res,
    point meta min=0,
  },
}


% ---- geometry, in points on the 510 x 200pt canvas ------------------------
\def\axL{34}      % x of the left column's axes
\def\axW{196}     % width of (a) and (b)
\def\axH{81.5}    % height of (a) and (b)
\def\axBb{26}     % y of (b), the lower left panel
\def\axBa{115.5}  % y of (a) = \axBb + \axH + 8pt gap
\def\cxL{256}     % x of (c) = \axL + \axW + 26pt
\def\cxW{220}     % width of (c)
\def\cxH{171}     % height of (c): bottom of (b) up to top of (a)

\newcommand{\utau}{$\langle\tau_{\rm err}\rangle\ [\tau_0]$}
\newcommand{\uq}{$\dot{Q}\ [k_{\rm B}T/\tau_0]$}

\pgfplotsset{
  fig5axis/.style={figaxis, anchor=south west, xmin=0.6, xmax=7.4,
                   xtick={2,4,6}},
  % marker on every production point, white-rimmed as in the original
  fig5series/.style={line width=1.4pt, mark=*, mark size=1.8pt,
                     mark options={draw=white, line width=0.4pt}},
}
\tikzset{
  fig5letter/.style={figpanellabel, anchor=north west, fill=white,
                     fill opacity=0.75, text opacity=1, rounded corners=1.5pt,
                     inner sep=1.8pt},
}

\begin{document}
\begin{tikzpicture}
\figcanvas{\figfullwidth}{200pt}

% =========================================================================
%  (a)  mean bit-flip time against chain length
% =========================================================================
\begin{axis}[fig5axis,
  at={(\axL pt,\axBa pt)}, width=\axW pt, height=\axH pt,
  ymode=log, ymin=1e2, ymax=3e9,
  ytick={1e2,1e5,1e8}, minor ytick={1e3,1e4,1e6,1e7,1e9},
  xticklabels={}, ylabel={\utau},
]
% the reaction-diffusion reference: the V_dd=1.4 slope between L=1 and L=3,
% continued straight, so the gap that opens above the data IS the sublinearity
\addplot[figgray, line width=1.1pt, dashed, forget plot]
  table[x=L,y=y] {fig5_ref.dat};
\node[figgray, font=\figannot, anchor=base] at (axis cs:6.35,3.0e7) {$\propto L$};
\addplot[fig5series, figvddA] table[x=L,y=t11] {fig5_tau.dat};
\addplot[fig5series, figvddB] table[x=L,y=t12] {fig5_tau.dat};
\addplot[fig5series, figvddC] table[x=L,y=t13] {fig5_tau.dat};
\addplot[fig5series, figvddD] table[x=L,y=t14] {fig5_tau.dat};
\node[fig5letter] at (axis description cs:0.045,0.955) {(a)};
\end{axis}

% =========================================================================
%  (b)  steady-state dissipation against chain length
% =========================================================================
\begin{axis}[fig5axis,
  at={(\axL pt,\axBb pt)}, width=\axW pt, height=\axH pt,
  ymin=0, ymax=150, ytick={0,50,100,150},
  xlabel={$L$}, ylabel={\uq},
]
\addplot[fig5series, figvddA] table[x=L,y=q11] {fig5_dissipation.dat};
\addplot[fig5series, figvddB] table[x=L,y=q12] {fig5_dissipation.dat};
\addplot[fig5series, figvddC] table[x=L,y=q13] {fig5_dissipation.dat};
\addplot[fig5series, figvddD] table[x=L,y=q14] {fig5_dissipation.dat};
\node[fig5letter] at (axis description cs:0.045,0.955) {(b)};
\end{axis}

% =========================================================================
%  (c)  the dissipation-reliability tradeoff
% =========================================================================
%  Limits are matplotlib's autoscale on this content: 5% margins, taken in
%  log space on the log axis (see the tail of prep_fig5.py).
\begin{axis}[figaxis, anchor=south west,
  at={(\cxL pt,\axBb pt)}, width=\cxW pt, height=\cxH pt,
  xmin=3.7138, xmax=139.9192, xtick={50,100},
  ymode=log, ymin=122.623, ymax=4.10104e6,
  ytick={1e3,1e4,1e5,1e6},
  yticklabel pos=right, ylabel near ticks,
  % a right-hand ylabel placed via `ticklabel cs. loses figaxis.s label font
  ylabel style={font=\figaxislabel},
  xlabel={\uq}, ylabel={\utau},
  legend style={at={(0.98,0.02)}, anchor=south east, cells={anchor=west},
                row sep=0.5pt},
  legend image post style={mark size=1.8pt},
]
\addplot[fig5series, line width=1.5pt, mark size=2pt, figvddA]
  table[x=q11,y=t11] {fig5_tradeoff.dat};
\addplot[fig5series, line width=1.5pt, mark size=2pt, figvddB]
  table[x=q12,y=t12] {fig5_tradeoff.dat};
\addplot[fig5series, line width=1.5pt, mark size=2pt, figvddC]
  table[x=q13,y=t13] {fig5_tradeoff.dat};
\addplot[fig5series, line width=1.5pt, mark size=2pt, figvddD]
  table[x=q14,y=t14] {fig5_tradeoff.dat};
\legend{$V_{\rm dd}/V_T=1.1$,$V_{\rm dd}/V_T=1.2$,%
        $V_{\rm dd}/V_T=1.3$,$V_{\rm dd}/V_T=1.4$}
\node[fig5letter] at (axis description cs:0.045,0.955) {(c)};
\end{axis}

\end{tikzpicture}
\end{document}
